\chapter{Conclusion}
\label{chap:conclusion}

This thesis investigated whether ZGC can process queue-less \texttt{WeakReference} objects more efficiently, and whether weak semantics expressed directly in object fields can reduce overhead further. To answer these questions, three orthogonal \texttt{WeakReference}-pipeline optimisations were implemented in an OpenJDK research fork: a queue-less fast path that skips enqueue handling, a dynamic-array representation for discovered queue-less weak references, and a specialised clear path that removes unnecessary per-reference work. In addition, an exploratory weak-field mechanism was implemented through a Java field annotation and corresponding VM support.

The evaluation shows a clear answer to \cref{rq:weakref-opt}. For the reference-processing subphase most directly targeted by the implementation, the strongest \texttt{WeakReference} variants are those that combine the specialised clear path with the dynamic array. In the single-object benchmark, \texttt{clear\_path\_dyn} and \texttt{all} reduce median non-strong processing time from 996.9\,ms to 187.9\,ms and 184.4\,ms respectively, corresponding to reductions of about 81\%. In the multi-object benchmark, the same variants reduce the median from 44.1\,ms to 18.8\,ms and 18.7\,ms, or about 57\%. The skip-enqueue separation on its own contributes little, which indicates that the main savings come from improved traversal locality and reduced per-reference work rather than from the list split alone.

For broader GC cycle metrics, the answer is more restrained. The best \texttt{WeakReference} variants reduce median major-collection time by approximately 8\% in the single-object benchmark, but those improvements are small relative to the overall spread of the measurements. In other words, optimising queue-less weak-reference processing clearly helps the targeted subphase, but it does not by itself transform the full cost profile of a GC cycle.

The answer to \cref{rq:weak-fields} is more mixed but also more structurally interesting. The weak-field mechanism does not achieve the lowest non-strong processing times; the best specialised \texttt{WeakReference} variants still perform better on that metric. However, weak fields substantially reduce Java heap usage and also improve total collection metrics, especially in the single-object benchmark. This indicates that a large share of weak-reference overhead comes from representation cost, namely the allocation and retention of millions of wrapper objects, rather than only from the mechanics of the non-strong processing phase.

Taken together, the results support two conclusions. First, ZGC's handling of queue-less \texttt{WeakReference} objects can be made materially faster with targeted, semantics-preserving implementation changes. Second, representation choices matter at least as much as pipeline mechanics: a field-based weak-semantics model can reduce broader GC costs by shrinking the heap structure that the collector has to traverse, even if its current prototype does not yet outperform the best specialised \texttt{WeakReference} path on every subphase. The thesis therefore contributes both a concrete optimisation study for ZGC's existing weak-reference pipeline and an exploratory argument that future JVM designs may benefit from supporting weak semantics more directly in object representation.

Placed in the broader context surveyed in \cref{chap:related-work}, these results also carry a wider implication. Across contemporary languages, callback-less weak semantics is the norm for the pure liveness-probing use case: Go's \texttt{weak.Pointer}~\cite{go-weak-pkg}, C++'s \texttt{std::weak\_ptr}~\cite{cppreference-weak-ptr}, .NET's \texttt{WeakReference<T>}~\cite{dotnet-weakreference}, and Lua's weak table entries~\cite{lua54-manual} all belong to this family, providing a probe mechanism without any notification callback. Java's \texttt{WeakReference} API already offers a callback-less path through its one-argument constructor, but the JVM has historically processed queue-less references identically to queue-backed ones, incurring the full cost of the notification machinery even when no notification was wanted. The queue-less optimisations developed here close that gap: ZGC can now handle references that were always conceptually equivalent to Go's \texttt{weak.Pointer} with correspondingly lower runtime overhead. The weak-field mechanism takes the convergence further, eliminating the wrapper object entirely and moving the representation closer to Lua's embedded weak entries or .NET's \texttt{ConditionalWeakTable}. Together, the two lines of work suggest a trajectory along which Java and ZGC can move towards the leaner weak-semantics designs found elsewhere in the language ecosystem, without abandoning the richer notification capabilities that distinguish Java's reference API for workloads that genuinely need them.
